\section{Ethical, Societal and Sustainability Considerations}\label{sec:ethics}
\textbf{Privacy and re-identification.} Health records are among the most sensitive data a state holds. The design limits risk at four points: identifiers are hashed with a secret salt at the bronze-silver boundary and exact age is discarded; the dashboard receives only aggregates; counts of one to four are suppressed with complementary suppression; and a district role cannot download other districts' detail because that detail is not in its files. Residual risks remain. A hashed identifier is still a pseudonym, so silver must stay inside the secure zone. Weekly district counts of five or more can still be linked with outside knowledge in very small districts, and repeated releases can be differenced. A public release would need formal protection such as differential privacy: adding calibrated noise to weekly counts (a Laplace mechanism with sensitivity one per person and a privacy budget split across the 261 weeks and 256 series, a large total) would make small-district series very noisy, which is why we do not claim it for public release and keep the dashboard restricted to authorised officials. A k-anonymity check~\cite{sweeney} on record-level extracts (PRV-4) was not implemented because the system never produces such extracts.

\textbf{Equity and misuse.} Allocation rules embed value judgements. Weighting unmet demand by hospital beds per capita directs scarce supplies to under-resourced districts, but an unconstrained optimiser can leave low-risk districts with nothing; the equity floor makes the trade-off explicit and the dashboard shows what each level of $\alpha$ costs in efficiency and what it buys in coverage of the worst-served district. Forecast error is not uniform: sparse districts (hill districts with few cases) have wide intervals, so a rule that favours high-count districts would systematically under-serve them. Alerts can stigmatise a district or community; the dashboard therefore shows alert status for officials only and avoids naming facilities. Models and thresholds could be misused to justify cutting support to districts with ``low'' forecasts; the policy brief states the limits of the forecasts.

\textbf{Sustainability and cost.} The platform runs on a laptop and a modest server, uses open-source software, and relies on columnar compression and partition pruning to avoid scanning unnecessary data. Precomputed scenario grids avoid running an optimisation server. A national deployment would need sustained funding for data quality at facility level, which matters more than model sophistication: the largest source of error in any surveillance system is missing or late reporting.

\textbf{AI assistance and authorship.} AI coding assistance was used and is documented in \texttt{AI\_USAGE.md}. Every member is responsible for being able to explain the code they present.
